\medskip\noindent\textbf{Exact correspondences for the two orbit-twenty detectors.}
The two orbit-twenty tests in (F.7.32) admit the following
explicit subgroup and filling description. It gives
continuous word substitutions and a common section of
the surface-group quotient maps. Faithfulness of the
outer-automorphism detectors remains unresolved.

Use the precise $E,q_5,H,K,a,b,d,e$ of (F.7.19)--(F.7.22).
Put $T=\langle q_5(b)\rangle$ and $B=q_5^{-1}(T)$.
The index-twenty orbifold cover has two order-three
points and four width-five cusps, and no order-two
points. This is the exact coset action: $a$ has ten
two-cycles, $b$ has two fixed cosets and six three-cycles,
and $d$ has four five-cycles. Its signature is
$(0;3,3;4\text{ cusps})$.

The following words give a specified full profinite
presentation:
\[
 \begin{split}
 z&=aba b^{-1},\qquad \tau=zaz^{-1},\\
 r&=b,\qquad s=\tau b\tau^{-1},\\
 u&=(ab^{-1})^4a,\qquad v=(ab)u(ba),\\
 w&=(aba b^{-1}a)(ba)^3(b^{-1}aba).
 \end{split}
 \qquad
 B=\langle r,s,u,v,w\mid r^3=s^3=1\rangle_{\rm prof}.
                                                               \tag{F.7.33}
\]
Thus $B=C_3\amalg C_3\amalg\widehat F_3$ with the
displayed inclusion into $E$. The exact Schreier
certificate uses twenty right cosets $Tg$ and the
prefix-closed breadth-first transversal in $a,b$.
Its forty directed generator edges have nineteen
tree edges, leaving twenty-one Schreier generators.
Rewriting $a^2,b^3$ at the twenty vertices gives
forty relators. Sixteen Tietze eliminations, each
solving a relator in a generator occurring exactly
once, leave the five displayed generators and
precisely the two cube relators.
The complete transversal, original generator words,
substitutions, and eliminations are recorded.
Reduction in $C_2*C_3$ checks the compact words in
(F.7.33). These are inverse presentation maps;
completion of the finite-index discrete inclusion
retains its full profinite topology.

The four based cusp loops and their actual
$T$-orbits of $H$ punctures are
\[
 \begin{array}{c|c|c}
 i&\kappa_i\text{ in }B&J_i\\ \hline
 0&u^{-1}r&\{0,1,2\}\\
 1&v^{-1}&\{3,4,6\}\\
 2&wu&\{5,7,8\}\\
 3&vw^{-1}s&\{9,10,11\}
 \end{array}                                                \tag{F.7.34}
\]
Each is the rewritten word $t_i d^5t_i^{-1}$
for its recorded transversal representative.
The first and fourth triples lie in the
$A_5$-orbit of $\{0,1,2\}$; the middle two lie
in the orbit of $\{0,4,9\}$.

Fill the other three orbifold cusps and put
$B_i=B/\langle\!\langle\kappa_j:j\ne i\rangle\!\rangle_B$.
To
$A=\langle R,S\mid R^3=S^3=1\rangle_{\rm prof}=C_3\amalg C_3$
send $r$ to $R$, $s$ to $S$, and use the indicated row:
\[
 \begin{array}{c|ccc|c}
 i&u&v&w&\kappa_i\\ \hline
 0&S^{-1}&1&S&SR\\
 1&R&S^{-1}R^{-1}&R^{-1}&RS\\
 2&R&1&S&SR\\
 3&R&1&R^{-1}&RS
 \end{array}                                                \tag{F.7.35}
\]
Substitution kills the three required cusp relators.
Conversely those three relators solve for $u,v,w$
as the entries of the row. Inclusion of $r,s$ in
the quotient gives the inverse map from $A$.
This proves the exact quotient isomorphism and
identifies the retained positive cusp as the
displayed product without a power ambiguity.

Let $E^+=\ker(\epsilon:E\to C_2)$, where
$\epsilon(a)=1,\epsilon(b)=0$.
The two-coset Schreier calculation with transversal
$1,a$ gives
\[
 A\xrightarrow{\ \sim\ }E^+,\qquad
                  R\longmapsto b,\quad S\longmapsto aba.
                                                               \tag{F.7.36}
\]
Write $p_i:B\to E^+$ for the composite of the row
(F.7.35) and (F.7.36).
The quotient $B/H=T$ sends
\[
 (r,s,u,v,w)\longmapsto
       (\bar b,\bar b^{-1},\bar b,1,\bar b^{-1}),
\]
whereas $E^+/K=C_3$ sends $b$ to its generator
and $aba$ to the inverse. Identify these two cyclic
groups by their named generators.
All four rows intertwine the specified quotient maps,
and hence $p_i(H)=K$.

The restriction to $H$ is precisely the filling
of the nine punctures outside $J_i$.
Indeed the closed $B$-normal closure of the other
three cusp words is contained in $H$.
Conjugation by $B/H=T$ supplies the three $H$ cusp
classes in each omitted orbit, and further
$H$ conjugation gives their closed normal closure.
Thus
\[
 \begin{split}
 \langle\!\langle\kappa_j:j\ne i\rangle\!\rangle_B
 &=\langle\!\langle c_\ell:\ell\notin J_i
                                      \rangle\!\rangle_H,\\
 1\longrightarrow H_{J_i}&\longrightarrow B_i
                 \longrightarrow C_3\longrightarrow1,
       \qquad H_{J_i}\xrightarrow[\ p_i\ ]{\ \sim\ }K .
 \end{split}                                                \tag{F.7.37}
\]
Each omitted cusp has width one for $H\leq B$:
$\kappa_j\in H$ and its three lifts have precisely
the labels in (F.7.34). No peripheral power is
inserted into this normal-closure identity.

There is a common actual section. Define
\[
 \iota_0:E\longrightarrow E,\qquad
                a\longmapsto\tau,\quad b\longmapsto b .
                                                               \tag{F.7.38}
\]
The orders give a homomorphism. On
$E^+=\langle b,aba\rangle$ it maps the two free
profinite factors to $r,s$, whose inclusion in
$B$ is injective by its retraction onto those factors.
Moreover $q_5(\tau)$ is an involution inverting
$q_5(b)$, so $q_5\iota_0(E)$ is the order-six
normalizer of $T$. Its quotient by $T$ is
exactly $\epsilon$. Thus $\ker\iota_0$ is contained
in $E^+$, where it is trivial, proving that
(F.7.38) is a continuous embedding.

The kernel of
$E\to q_5\iota_0(E)\cong S_3$ is the characteristic
$K$ proved in (F.7.20). Hence $\iota_0(K)\leq H$.
Every $p_i$ is the identity on the displayed two
factors after (F.7.36), so
\[
 p_i\iota_0|_{E^+}=1_{E^+},\qquad
             p_i|_H\circ\iota_0|_K=1_K.                 \tag{F.7.39}
\]
The program also rewrites $\iota_0(x),\iota_0(y)$
through all four quotient rows and verifies that
their images are exactly $x,y$.

For $\gamma=(1,f)\in GT_0^1$, the actual normalized
automorphism (F.7.27) extends to an automorphism of
$E$ inducing the identity on $A_5$. It preserves
$B$ and each filled normal subgroup.
Using (F.7.37) for the target identification, the
actual detector is
\[
 \mathcal D_{J_i}(\gamma)
        =\bigl[p_i\,\eta_f^H\,\iota_0|_K\bigr]
                          \quad\hbox{in }\mathcal P_K^1.
                                                               \tag{F.7.40}
\]
Here $\eta_f^H$ is applied to $\iota_0(K)\leq H$.
Its preservation of that embedded subgroup is
not assumed. The subsequent quotient map is
part of the composition. The composition is
an automorphism because it is the actual descended
automorphism under (F.7.37), evaluated on the
section (F.7.39).
The two pairs of equal detector kernels are
the actual orbit transports of (F.7.31)--(F.7.32).

The four source quotient maps themselves have
a nontrivial common kernel:
\[
 [\kappa_0,\kappa_1]
           =u^{-1}r v^{-1}r^{-1}uv\ne1.                 \tag{F.7.41}
\]
Every row kills this word.
Its nontriviality in the full completion is
detected by $B\to S_3$ sending $r,s,w$ to $1$,
$u$ to $(1\ 2)$, and $v$ to $(2\ 3)$:
its image has order three.
This source word gives an inner automorphism
of $E$, whose class in $\operatorname{Out}(E)$
is the identity. It therefore supplies no
nonidentity $GT_0^1$ element in the kernel of
(F.7.40). Faithfulness of an outer detector,
or of the five-detector family, remains open.

The executed program
\path{code/part_a/modular_c3_correspondence_f7.py}
records the full Schreier presentation, its
sixteen eliminations, the four three-step filling
eliminations, compact word identities, exact
finite quotient labels, eight common-section
word checks, and the finite detector in (F.7.41).
